\begin{table}[H]\centering\footnotesize
\caption{Scaling on MNIST-10k (Scenario 2).}
\label{tab:scale}
\begin{tabular}{r rrr c}
\toprule
& \multicolumn{3}{c}{wall-clock time} & accuracy\\
\cmidrule(lr){2-4}
$N$ & dense \texttt{eigh} & \texttt{svds} & neural & AMI (svds\,$\approx$\,neural)\\
\midrule
2\,000  & 4.3\,s  & 0.07\,s & 18.5\,s & 0.74\\
4\,000  & 14.6\,s & 0.11\,s & 40.0\,s & 0.79\\
8\,000  & 98.0\,s & 0.30\,s & 71.3\,s & 0.81\\
10\,000 & n/a     & 0.31\,s & 89.4\,s & 0.81\\
\bottomrule
\end{tabular}
\end{table}

\noindent\textbf{Setup.} Time to build the embedding of $N$ points three ways, plus the resulting
$k$-means AMI. dense \texttt{eigh} = classical $O(N^3)$ eigendecomposition of the dense operator;
\texttt{svds} = truncated top-$k$ SVD of the sparse, matrix-free operator action (Scenario 2);
neural = the Rayleigh-trained net (Scenario 1). \texttt{svds} and neural produce the same
embedding, hence one shared accuracy column; dense \texttt{eigh} is omitted at $N{=}10$k (too slow).

\noindent\textbf{Result.} \texttt{svds} ($0.07$--$0.31$\,s) is ${\sim}300\times$ faster than dense
\texttt{eigh} and ${\sim}200\times$ faster than the net, at identical accuracy ($0.74$--$0.81$). The
scalable classical method wins; the neural machinery buys no speed.
